% !TEX root = ../thesis.tex
% ==============================================================================
%  APPENDIX: CONTRASTIVE NEGATIVE SAMPLING STRATEGIES (appendices/appendix_negative_samplers.tex)
% ==============================================================================
%  This appendix details the contrastive negative sampling mechanisms used
%  to train the energy-based surrogate model (EBM) via InfoNCE loss.
% ==============================================================================

\section{استراتژی‌های نمونه‌برداری تقابلی در آموزش شبکه انرژی}\label{app:appendix_negative_samplers}

آموزش صحیح مدل مبتنی بر انرژی (\lr{EBM}) وابستگی عمیقی به مجموعه‌ی نمونه‌های منفی استفاده‌شده در تابع زیان \lr{InfoNCE} دارد. برای دستیابی به خمینه‌ی انرژی مطلوب، ما مجموعه‌ای متنوع از استراتژی‌های نمونه‌برداری طراحی کرده‌ایم که هر یک ویژگی خاصی از پویاییِ تولیدی و استنتاجی مدل زبانی را هدف قرار می‌دهند. این نمونه‌بردارها به دو دسته‌ی کلیِ «مبتنی بر مجموعه‌داده» (استفاده از نمونه‌های خارجی یا سایر متون درون بچ) و «مبتنی بر اختلال رویه‌ای» (اعمال تغییرات ساختاری بر روی خود متن) تقسیم می‌شوند. 

در آزمایش‌های این پایان‌نامه، تأثیر سه پیکربندی مجزا از این نمونه‌بردارها مورد بررسی قرار گرفته است. در ادامه، مکانیزم‌های استفاده‌شده در این پژوهش به‌طور دقیق و با ارائه‌ی فرمول‌بندی‌های ریاضی تشریح می‌گردند. فرض کنید بچ آموزشیِ ما دارای اندازه‌ی $B$ بوده و شامل جفت‌های اصیل $(\mathbf{x}_i, \mathbf{y}_i)$ برای $i \in \{1, \dots, B\}$ است.


\subsection{نمونه‌بردارهای پیکربندی وفاداری}

این مجموعه ($\mathcal{M}_{\text{\lr{Fidelity}}}$) بر آموزش توزیع کلان مدل تمرکز دارد و تلاش می‌کند شبکه‌ی جانشین را قادر سازد تا امضای تولیدی و سبک اختصاصی مدل هدف را شناسایی کند. این پیکربندی شامل مکانیزم‌های زیر است:

\begin{itemize}
    \item \textbf{جایگزینی با پاسخ انسانی (\texttt{response\_human}):} 
    در این روش، پرامپت $\mathbf{x}_i$ ثابت نگه داشته می‌شود و پاسخ تولیدشده توسط مدل جعبه‌سیاه ($\mathbf{y}_i$) با پاسخ انسانی متناظر در مجموعه‌داده‌ی مرجع ($\mathbf{y}^{\text{\lr{hum}}}_i$) جایگزین می‌گردد. این نمونه‌بردار به مدل کمک می‌کند تا تفاوت‌های ظریف میان متن تولیدشده توسط انسان و مدل را فراگیرد:
    \begin{equation}
        \mathcal{N}^{\text{\lr{hum}}}_i = \{ (\mathbf{x}_i, \mathbf{y}^{\text{\lr{hum}}}_i) \}
    \end{equation}

    \item \textbf{جایگزینی با پاسخ مدل‌های دیگر (\texttt{response\_other\_lm}):}
    مشابه روش قبل، پرامپت $\mathbf{x}_i$ ثابت است، اما این بار پاسخ با خروجی ازپیش‌محاسبه‌شده‌ی یک مدل زبانی متفاوت (نظیر \lr{GPT-2}) برای همان پرامپت ($\mathbf{y}^{\text{\lr{other}}}_i$) جایگزین می‌شود تا \lr{EBM} لحن اختصاصی مدل هدف را بیاموزد:
    \begin{equation}
        \mathcal{N}^{\text{\lr{other}}}_i = \{ (\mathbf{x}_i, \mathbf{y}^{\text{\lr{other}}}_i) \}
    \end{equation}

    \item \textbf{اختلال پرامپت و پاسخ درون‌بچ (\texttt{off\_topic}):}
    در این روش، هدف از بین بردن ارتباط موضوعی میان پرامپت و پاسخ است. برای هر جفت اصیل، یکی از مؤلفه‌ها ثابت می‌ماند و مؤلفه‌ی دیگر با متون متناظر از تمامیِ $B-1$ نمونه‌ی دیگرِ درون همان بچ جابه‌جا می‌شود. این استراتژی دو مجموعه‌ی منفیِ متمایز تولید می‌کند؛ نخست، با ثابت نگه‌داشتن پرامپت، پاسخ‌ها جابه‌جا می‌شوند (\texttt{response\_off\_topic}):
    \begin{equation}
        \mathcal{N}^{\text{\lr{off-topic-resp}}}_i = \left\{ (\mathbf{x}_i, \mathbf{y}_j) \;\middle|\; j \in \{1, \dots, B\}, j \neq i \right\}
    \end{equation}
    و دوم، نسخه‌ی متناظر آن برای پرامپت (\texttt{prompt\_off\_topic}) که با ثابت نگه‌داشتن پاسخ، پرامپت‌ها را جابه‌جا می‌کند:
    \begin{equation}
        \mathcal{N}^{\text{\lr{off-topic-pmt}}}_i = \left\{ (\mathbf{x}_j, \mathbf{y}_i) \;\middle|\; j \in \{1, \dots, B\}, j \neq i \right\}
    \end{equation}

    \item \textbf{ماسک‌گذاری جمله‌ای (\texttt{sentence\_masking}):}
    در این رویکرد رویه‌ای، ساختار خود متن مخدوش می‌شود تا مدل وابستگی به پیوستگی گرامری را بیاموزد. در این روش، هر جمله‌ی متن به‌صورت مستقل با احتمال $p_{\text{\lr{drop}}}$ (به‌طور پیش‌فرض $0.3$) حذف می‌گردد تا یک متن ناقص ایجاد شود. الگوریتم تضمین می‌کند که همواره حداقل یک جمله حذف شده و حداقل یک جمله در متنِ نهایی باقی بماند تا از فروپاشی کامل ورودی جلوگیری شود. این استراتژی بر روی پاسخ‌ها (\texttt{response\_sentence\_masking}) و پرامپت‌ها (\texttt{prompt\_sentence\_masking}) به‌طور مجزا اعمال می‌شود:
    \begin{align}
        \mathcal{N}^{\text{\lr{mask-resp}}}_i &= \{ (\mathbf{x}_i, \mathbf{y}^{\text{\lr{mask}}}_i) \} \\
        \mathcal{N}^{\text{\lr{mask-pmt}}}_i &= \{ (\mathbf{x}^{\text{\lr{mask}}}_i, \mathbf{y}_i) \}
    \end{align}
\end{itemize}


\subsection{نمونه‌بردارهای پیکربندی وابستگی محلی}

این مجموعه ($\mathcal{M}_{\text{\lr{Dep}}}$)، چنانچه در بخش~\ref{sec:4-3-2} نیز مورد بحث قرار گرفت، مشخصاً برای کشف پیش‌نیازهای علی و روابط منطقی در سطح جزء‌به‌جزء طراحی شده است. این پیکربندی از ترکیب مکانیزم ماسک‌گذاری تصادفی و تقابل درون‌بچ برای تولید نمونه‌های پادواقعی بهره می‌برد:

\begin{itemize}
    \item \textbf{تقابل وابستگی پاسخ جزئی (\texttt{partial\_response\_dep}):}
    ابتدا در طی یک مرحله‌ی پیش‌پردازش، تمامی پرامپت‌ها و پاسخ‌های درون بچ تحت عملیات ماسک‌گذاری تصادفی قرار می‌گیرند تا پرامپت‌های جزئی ($\mathbf{x}_{\text{\lr{part}}}$) و پاسخ‌های جزئی ($\mathbf{y}_{\text{\lr{part}}}$) شکل گیرند. سپس برای نمونه‌ی $i$، پرامپت جزئی آن ($\mathbf{x}_{\text{\lr{part}}, i}$) ثابت نگه داشته شده و در تقابل با پاسخ‌های جزئیِ سایر نمونه‌های درون بچ ($\mathbf{y}'_{\text{\lr{part}}, j}$) قرار می‌گیرد. این مکانیزم برای هر نمونه، $B-1$ نمونه‌ی منفی تولید می‌کند:
    \begin{equation}
        \mathcal{N}^{\text{\lr{resp-dep}}}_i = \left\{ \left(\mathbf{x}_{\text{\lr{part}}, i}, \mathbf{y}'_{\text{\lr{part}}, j}\right) \;\middle|\; j \in \{1, \dots, B\}, j \neq i \right\}
    \end{equation}
    این تقابل، شبکه را مجبور می‌کند تا ارزیابی نماید که آیا خروجی ارائه‌شده، واقعاً ادامه‌ی منطقیِ همان مفاهیمِ محدود در ورودی است یا صرفاً شباهت ظاهری دارد.

    \item \textbf{تقابل وابستگی پرامپت جزئی (\texttt{partial\_prompt\_dep}):}
    در این مکانیزم که آینه‌ی روش پیشین است، پاسخ جزئیِ نمونه‌ی $i$ام ($\mathbf{y}_{\text{\lr{part}}, i}$) ثابت می‌ماند و در برابر پرامپت‌های جزئیِ سایر نمونه‌های بچ ($\mathbf{x}'_{\text{\lr{part}}, j}$) قرار می‌گیرد:
    \begin{equation}
        \mathcal{N}^{\text{\lr{pmt-dep}}}_i = \left\{ \left(\mathbf{x}'_{\text{\lr{part}}, j}, \mathbf{y}_{\text{\lr{part}}, i}\right) \;\middle|\; j \in \{1, \dots, B\}, j \neq i \right\}
    \end{equation}
    این عملیات تضمین می‌کند که شبکه، پاسخ را منحصراً به پیش‌نیازهای علی واقعی آن در پرامپت انتساب دهد، نه اینکه به شباهت‌های موضوعی کلان در سطح بچ بسنده کند.
\end{itemize}


\subsection{نمونه‌بردارهای پیکربندی ترکیبی}

این پیکربندی ($\mathcal{N}_{\text{\lr{hybrid}}}$) که در نهایت به عنوان حالت پایدار و منتخب در این پژوهش معرفی شده است، ترکیبی انطباقی از رویکرد وفاداری و رویکرد وابستگی محلی است. شایان ذکر است که در این پیکربندی نهایی، ما از گنجاندن نمونه‌بردارهای خامِ \texttt{off\_topic} و \texttt{sentence\_masking} (که متعلق به فاز وفاداری بودند) صرف‌نظر کرده‌ایم. دلیل این امر آن است که اثرات اختلالات و پرش‌های موضوعی این دو روش، به‌طور ضمنی و هم‌زمان در دل ساختار پادواقعی نمونه‌بردارهای وابستگی محلی ادغام شده است. استفاده‌ی مجدد و مجزا از آن‌ها، صرفاً به پیچیدگی تابع زیان و ایجاد نویز غیرضروری در فرآیند یادگیری می‌انجامد.

بر این اساس، در این حالت، برای هر نمونه‌ی $i$ام، مجموعه‌ی منفی نهایی منحصراً از اجتماع مستقیم نمونه‌های پادواقعی و نمونه‌های خارجی تشکیل می‌شود:
\begin{equation}
    \mathcal{N}^{\text{\lr{hybrid}}}_i = \mathcal{N}^{\text{\lr{resp-dep}}}_i \cup \mathcal{N}^{\text{\lr{pmt-dep}}}_i \cup \mathcal{N}^{\text{\lr{hum}}}_i \cup \mathcal{N}^{\text{\lr{other}}}_i
\end{equation}
ترکیب این استراتژی‌ها، ضمن حفظ تعادل میان یادگیری ساختار توزیعی کلان مدل هدف و کشف دقیق روابط علی، \lr{EBM} را به یک ارزیاب مستحکم و قابل‌اتکا تبدیل می‌نماید.
